\section{SDD Template --- Nogoolin}\label{doc6-sdd-template-nogoolin}

\textbf{ICSI405 / M2 · Тэнгис · 2026-09-15 · v0.1 / draft}

\subsection{1. Зорилго ба шаардлагын
холбоо}\label{doc6-ux437ux43eux440ux438ux43bux433ux43e-ux431ux430-ux448ux430ux430ux440ux434ux43bux430ux433ux44bux43d-ux445ux43eux43bux431ux43eux43e}

Энэ SDD нь \href{../../requirements/requirements.md}{таван шаардлагыг}
хэрэгжүүлэх бүтэц, технологийг тайлбарлана. M2-д одоогийн кодоос үндсэн
бүтэц, design decision, gap-ийг тогтоов; arc42-ийн дэлгэрэнгүй W4,
OpenAPI W5, диаграм W8-д төлөвлөгдсөн. SRS нь хэрэглэгчийн авах үр дүн,
SDD нь түүнийг хэрэгжүүлэх арга.

\subsection{2. Системийн орчин ба
бүрэлдэхүүн}\label{doc6-ux441ux438ux441ux442ux435ux43cux438ux439ux43d-ux43eux440ux447ux438ux43d-ux431ux430-ux431ux4afux440ux44dux43bux434ux44dux445ux4afux4afux43d}

\begin{longtable}[]{@{}
  >{\raggedright\arraybackslash}p{(\linewidth - 4\tabcolsep) * \real{0.2600}}
  >{\raggedright\arraybackslash}p{(\linewidth - 4\tabcolsep) * \real{0.5500}}
  >{\raggedright\arraybackslash}p{(\linewidth - 4\tabcolsep) * \real{0.1900}}@{}}
\toprule\noalign{}
\begin{minipage}[b]{\linewidth}\raggedright
Бүрэлдэхүүн
\end{minipage} & \begin{minipage}[b]{\linewidth}\raggedright
Шийдвэр / бодит байршил
\end{minipage} & \begin{minipage}[b]{\linewidth}\raggedright
Холбогдох шаардлага
\end{minipage} \\
\midrule\noalign{}
\endhead
\bottomrule\noalign{}
\endlastfoot
Web + admin & Next.js/TypeScript; \nolinkurl{apps/web}; inquiry,
wishlist, profile, admin inbox & FR-01/02, NFR-01/02 \\
API & Fastify/TypeScript; \nolinkurl{backend/api}; \nolinkurl{/api/v1},
local port 3001 & FR-01/02, NFR-01 \\
Өгөгдөл/Auth & Supabase local PostgreSQL/Auth/Storage;
\nolinkurl{supabase/migrations} & FR-01/02, NFR-01 \\
Shared validation & Zod; \nolinkurl{packages/validation-schemas} &
FR-01/02 \\
Hero & Three.js/R3F; \nolinkurl{apps/web/src/components/intro} &
NFR-02 \\
Build/CI & pnpm workspaces, Docker,
\nolinkurl{.github/workflows/api-deploy.yml} & NFR-03 \\
\end{longtable}

\textbf{Шийдвэрийн эх:}
\href{../../../agent-context/DECISIONS.md}{ADR-003/004/005/006/007/008/011}.
Mobile нь React Native + Expo prebuild/dev client; энэ M2 багц mobile
inquiry/wishlist-ийг дууссан гэж үзэхгүй. Railway/Vercel/cloud Supabase
нь Phase 6-ын байршуулалтын зорилт, одоо ажиллаж буй production орчин
гэж батлаагүй.

\subsection{3. Backend-ийн дотоод
бүтэц}\label{doc6-backend-ux438ux439ux43d-ux434ux43eux442ux43eux43eux434-ux431ux4afux442ux44dux446}

Controller → Service → Repository. Controller нь request/response ба
validation hook; service нь бизнесийн дүрэм; repository нь өгөгдөлд
хүрэх кодыг хариуцна. Auth hook хэрэглэгчийн session-ийг тогтооно.
Бизнесийн логикийг route руу, DB query-г service рүү холихгүй.

\textbf{FR-01:} inquiry controller → inquiry service → inquiry
repository → \nolinkurl{inquiries}. \textbf{FR-02:} cart controller →
cart service → cart repository → \nolinkurl{cart_items}. Shared schema-г
web/API нэг эхээс хэрэглэнэ; client validation нь server validation-ийг
орлохгүй.

\subsection{4. Өгөгдөл ба эрхийн
дизайн}\label{doc6-ux4e9ux433ux4e9ux433ux434ux4e9ux43b-ux431ux430-ux44dux440ux445ux438ux439ux43d-ux434ux438ux437ux430ux439ux43d}

\nolinkurl{inquiries.customer_id} нь нэвтэрсэн илгээгчийг заана, guest-д
null; шинэ асуулгын төлөв \nolinkurl{new}. \nolinkurl{cart_items} нь
\nolinkurl{(user_id, product_id)} unique хостой. Нэмэх үйлдлийг давтахад
нэг бичлэг хэвээр; хасах нь давтаж дуудаж болох үйлдэл. Эх:
\href{../../../supabase/migrations/20260726120000_inquiry_customer_and_cart.sql}{20260726120000
migration},
\href{../../../backend/api/src/repositories/supabase/cart.repository.ts}{cart
repository},
\href{../../../backend/api/src/repositories/supabase/inquiry.repository.ts}{inquiry
repository}.

\textbf{NFR-01:} customer ID-г request body-оос авахгүй, баталгаажсан
session-ээс дамжуулна. Cart read/delete нь
\nolinkurl{.eq('user_id', userId)}, inquiry mine нь
\nolinkurl{.eq('customer_id', customerId)}; insert session-ийн ID
онооно. RLS нь DB түвшний хамгаалалт, admin scope нь RBAC шалгалттай.
Бүх query-д нэг ижил \nolinkurl{user_id} шүүлтүүр шаардахгүй: public
catalog, guest create, admin ажиллагаа тус тусын эрхийн дүрэмтэй.

\clearpage
\subsection{5. Гадаад интерфэйс ба contract
gap}\label{doc6-ux433ux430ux434ux430ux430ux434-ux438ux43dux442ux435ux440ux444ux44dux439ux441-ux431ux430-contract-gap}

\href{../../phase-0/06-api-spec.yaml}{Phase 0 OpenAPI} нь эх contract.
Код дахь \nolinkurl{/cart} болон \nolinkurl{/inquiries/mine}, customer
ownership болон admin filter нэмэлтүүдийг W5-д тулган шинэчилнэ. Inquiry
нь email field-гүй; UI wireframe-ийн хуучин email мөрийг хэрэгжилттэй
ижил гэж үзэхгүй. \nolinkurl{INQ-YYYY-NNNN} нь UI-ийн derived display
утга бөгөөд дараалсан захиалгын дугаар биш.

\subsection{6. Hero fallback-ийн
дизайн}\label{doc6-hero-fallback-ux438ux439ux43d-ux434ux438ux437ux430ux439ux43d}

GLB URL тохируулаагүй үед \nolinkurl{deity.tsx} геометр орлуулагч
сонгоно. \nolinkurl{hero-intro.tsx} нь WebGL/reduced-motion-оор static
home төлөв сонгож, skip/catalog холбоосыг DOM дээр харуулна. Бодит
asset-ийг \href{../../../agent-context/ASSET_SPECS.md}{ASSET\_SPECS}-ийн
дагуу дараа солино. Алдаатай GLB URL-ийн error boundary/recovery-г энэ
draft баталгаажаагүй gap гэж тэмдэглэв. Product 360° viewer нь D-07-оор
post-MVP; нүүрний intro-оос тусдаа.

\subsection{7. Build, deploy ба
verification}\label{doc6-build-deploy-ux431ux430-verification}

API Dockerfile нь multi-stage build, non-root runtime ашигладаг.
Одоогийн CI dependency нь
\nolinkurl{lint-and-test → docker-build → deploy}.
\href{../../../backend/api/tests/inquiry-cart.integration.test.ts}{Тестийн
эх} DB хүрэхгүй үед skip хийдэг; workflow local Supabase бэлтгэхгүй тул
NFR-03-ын all-skipped хаалт хангагдаагүй. W9-д CI DB setup,
skip-fails-CI нөхцөл, гурван сөрөг шалгалтын log төлөвлөсөн. Production
auto-deploy нь энэ хаалтыг тойрохгүй байх тохиргоог Phase 6-д шалгана.

Энэ task-аар код, DB migration, CI-ийн ажиллагааг өөрчлөөгүй; шинэ
runtime/build тест ажиллуулаагүй.
\href{../../../agent-context/PROGRESS.md}{PROGRESS}-ийн өмнөх тестийн
тайлан нь түүхэн тэмдэглэл, одоогийн ажиллуулалтын нотолгоо биш.

\subsection{8. Үргэлжлүүлэн бөглөх
хэсгүүд}\label{doc6-ux4afux440ux433ux44dux43bux436ux43bux4afux4afux43bux44dux43d-ux431ux4e9ux433ux43bux4e9ux445-ux445ux44dux441ux433ux4afux4afux434}

\begin{longtable}[]{@{}
  >{\raggedright\arraybackslash}p{(\linewidth - 4\tabcolsep) * \real{0.2600}}
  >{\raggedright\arraybackslash}p{(\linewidth - 4\tabcolsep) * \real{0.5500}}
  >{\raggedright\arraybackslash}p{(\linewidth - 4\tabcolsep) * \real{0.1900}}@{}}
\toprule\noalign{}
\begin{minipage}[b]{\linewidth}\raggedright
Хэсэг
\end{minipage} & \begin{minipage}[b]{\linewidth}\raggedright
Нөхөх агуулга
\end{minipage} & \begin{minipage}[b]{\linewidth}\raggedright
Status
\end{minipage} \\
\midrule\noalign{}
\endhead
\bottomrule\noalign{}
\endlastfoot
Architecture views & arc42 context, building blocks, runtime,
deployment; FR/NFR холбоо & planned (W04) \\
API contracts & cart/history schema ба бодит error/status нийцүүлэлт &
planned (W05) \\
Diagrams & C4, ER, inquiry/auth sequence source ба review & planned
(W08) \\
Design verification & CI negative-case evidence, Docker build evidence,
traceability шинэчлэлт & planned (W09) \\
\end{longtable}

\textbf{Эх:} \href{../../phase-0/README.md}{Phase 0 index},
\href{../../../agent-context/DECISIONS.md}{шийдвэрүүд},
\href{../../requirements/srs-template.md}{SRS},
\href{../../requirements/traceability-matrix.md}{traceability}.
